\chapter{Literature Review}

% The literature review is a survey of the history and state of the art in the domain of your project. It will summarize the work that has already been done in the field; this may be scientific literature, known techniques, and even previous student projects. It will provide a historical perspective on how the subject area has arrived at its current state by looking at important developments over time. If appropriate, it may examine existing software in the domain, especially in terms of the technology used and the features offered. The focus of the literature review is to summarise the existing arguments and ideas of others, identifying which are important.

% A good literature review could be a project on its own, and form a very useful guide to anyone new to the particular field. It would identify the important work, authors and publications which would be a good place to begin research activity. Open questions and areas where new work is required would be discussed. Really good reviews are often published in scientific journals. Your literature review is not expected to be quite so substantial, but should still provide a comprehensive summary which will allow the reader to understand the field.

\section{Review}
Analysis of existing dynamic software updating techniques for safe and secure industrial control systems

Challenges in Dynamic Software Updating

A survey of dynamic software updating

\section{Case Study}

\subsection{Control Systems}
Live Seamless Firmware Upgrade in Real Time Control Systems 

\subsection{Dynamic Updates}
Dynamic Software Updates for Real-Time Systems.

\subsection{Disruption-free}
Disruption-free software updates in automation systems

\subsection{A model}
A Model for Updating Real-Time Applications

\subsection{Grid-connected}
Enhanced Uptime and Firmware Cybersecurity for Grid-Connected Power Electronics

\section{Rational}